\section{Выводы}
В ходе работы я познакомился с персистентными структурами данных и реализовал персистентное дерево отрезков.
Главное, что я понял: чтобы сохранить доступ ко всем прошлым версиям структуры, не нужно копировать её целиком.
При изменении достаточно скопировать путь от корня до изменяемого листа ($O(\log n)$ вершин), а остальные поддеревья
использовать совместно со старой версией. Этот приём называется копированием пути.

Я научился сводить геометрическую задачу (сколько отрезков лежит над точкой) к запросам на версиях дерева с
помощью метода заметающей прямой: события открытия и закрытия отрезков превращаются в версии, а запрос по
координате $x$ --- в выбор нужной версии бинарным поиском. Заодно я разобрал типичную ошибку с граничными случаями: если
закрывать отрезок в точке $r$, то при равных координатах событий результат зависит от порядка сортировки, а
закрытие в точке $r + 1$ (полуоткрытый отрезок) делает решение независимым от этого порядка.

Отдельно я разобрался с управлением памятью в персистентных структурах. Вершины общие у нескольких версий,
поэтому рекурсивное удаление от каждого корня приводит к повторному освобождению, а \texttt{shared\_ptr} слишком
расточителен по памяти и времени. Для задачи, где все версии живут до конца программы, лучше всего подошла арена:
все вершины лежат в одном \texttt{std::vector}, а ссылки между ними являются индексами. Это сократило расход
памяти (вершина занимает $12$ байт вместо $24$) и избавило от ручного освобождения. При этом пришлось учесть,
что при перевыделении вектора ссылки на его элементы становятся недействительными.

Практическое сравнение показало разницу между асимптотиками $O(nm)$ и $O((n + m)\log n)$: на максимальных
данных решение работает за доли секунды, а наивный перебор --- почти минуту. Подобные структуры пригодятся в
задачах с запросами к прошлым состояниям данных (запросы вида \enquote{на момент времени $t$}), в онлайн-вариантах
задач на подсчёт в диапазонах, а также в системах контроля версий и
баз данных, где неизменяемые структуры с общими частями применяются очень широко.

\pagebreak
